\subsection{Projection depth for separated windows}\label{dep:section}

We retain the notation and contradiction regime \eqref{red:regime}.
Let $P_x,P_y$ be the two pieces selected by
Proposition~\ref{prof:pair-selection}, in separated windows of widths
$v$ and $s$, respectively. Thus $v,s\to\infty$, both are
$o(\log B)$, and either $v/s\to0$ or $v/s\to\infty$.
The limiting normalized drop profile is the same for both pieces.
Before using the target plateaux, omit from $P_y$ the points left
unassigned by the global terminal assignment in
Proposition~\ref{prof:pair-selection}. Its proof bounds their number
by $o(|P_y|)$ uniformly over the pieces, so this removes
$o(|P_x||P_y|)$ ordered pairs. Continue to write $P_y$ for the
retained target set. It has size $N\exp(-o(s))$ and inherits the
profile conclusions, since it is a $(1-o(1))$ fraction of the
selected piece. The proposition below uses this retained $P_y$.

For the target piece write
\[
 D_i=B\exp(b_i s),\quad e_j=\prod_{i\le j}D_i,
 \quad l=d-j,
\]
and consider a positive-dimensional plateau after at least one drop.
The path equations defining its current system have degrees
$B\exp(O(s))$. For this plateau, fix one containing
$l$-dimensional carrier for each point of $P_y$. On the final
one-dimensional plateau use its carrier from the global terminal
assignment. Pool repeated carriers and denote the resulting family
by $\mathcal Y$, with disjoint assigned target sets of masses $n_Y$.
Thus $\sum_{Y\in\mathcal Y}n_Y=|P_y|$. All later incidence and list
restrictions use this fixed assignment.
Our objective is a lower bound for the degree of the distinct
projected components carrying any fixed positive fraction of target
mass.  Several components may share a projected image, and each
projection may have degree greater than one.  We control this combined
multiplicity after deleting negligible incidence mass.

Give a center--component incidence $(p,Y)$ normalized weight
$n_Y/(|P_x||P_y|)$.  The weight of a collection is thus the probability
that a uniformly chosen center and a uniformly chosen assigned target
point determine an incidence in that collection.

We use the following precise interpretation of exponential accuracy.
An assertion with a factor $\exp(o(s))$ in a lower bound means that
for every fixed $\epsilon>0$ it holds with the smaller factor
$\exp(-\epsilon s)$ for all sufficiently large members of the
sequence, after the stated vanishing-mass deletions.
All finite choices of plateaux, density constants, and positive
logarithmic tolerances can be imposed simultaneously.

\begin{proposition}[Projected degree in separated windows]
\label{dep:projected-degree}
For each fixed positive tolerance $\epsilon$, one can discard
center--target-component incidences of total normalized weight
$o(1)$, and then discard $o(|P_x|)$ centers, so that the following
holds. For each remaining center $p$, every target component $Y$
whose incidence $(p,Y)$ remains satisfies
$\dim\overline{\pi_pY}=l$. For every subcollection of those components
carrying at least a fixed positive fraction of $|P_y|$, the distinct
projected components have summed degree at least
\begin{equation}\label{dep:degree-bound}
 \frac{e_j}{R}\exp(-\epsilon s).
\end{equation}
The assertion is uniform over those subcollections, including ones
chosen after $p$ and after subsequent equations have been chosen.
The following values of $R$ are valid:
\begin{align}
 R&=D_1 &&\text{in either orientation};
       \label{dep:first-R}\\
 R&=B\exp\left(-s\sum_{i>j}b_i\right)
       &&\text{if }v/s\longrightarrow0;
       \label{dep:narrow-R}\\
 R&=B\exp(-2s)
       &&\text{if }v/s\longrightarrow\infty,
          \quad b_{j+1}>0,\quad l<d-1.
       \label{dep:wide-R}
\end{align}
In \eqref{dep:degree-bound}, ``fixed positive fraction'' means any
constant $\rho>0$ independent of the sequence. After fixing
$\epsilon$ and any finite batch of plateaux and tolerances, the
incidence and center deletions are independent of the later fixed
choice of $\rho$; only the large-index threshold may depend on
$\rho$.
\end{proposition}

\begin{proof}
We first explain the degree and mass bookkeeping common to all three
cases, then prove their depth bounds.

\paragraph{Components, multiplicity, and incidence deletions.}
Proposition~\ref{prof:budget} gives
\begin{equation}\label{dep:carrier-budget}
 \sum_{Y\in\mathcal Y}\deg Y\le e_j\exp(o(s)),
 \qquad
 \sum_{Y\in\mathcal A} n_Y
 \le \deg(\mathcal A)\left(\prod_{i>j}D_i\right)\exp(o(s))
\end{equation}
for every subcollection $\mathcal A$.
Since $\sum_i b_i=0$ and $|P_y|=N\exp(-o(s))$, every subcollection
of fixed positive assigned mass has degree at least
$e_j\exp(-o(s))$.

Every $Y\in\mathcal Y$ is an irreducible component of the fixed target
path locus, by the component argument at the start of the proof of
Proposition~\ref{prof:regularity}. Larger components elsewhere remain
in the target equations.

Projection from $p$ can drop the dimension of $Y$ only if its generic
fiber is a line. For a fixed general smooth $y\in Y$, this puts $p$
in the proper affine tangent flat $y+T_yY$. The flat cap makes the
proportion of such centers at most $B^{-c_0+o(1)}$, uniformly for each
$Y$. Also, Corollary~\ref{poly:high-multiplicity} puts centers where
$\operatorname{mult}_pY>\deg(Y)/2$ in a hypersurface of bounded degree.
The initial full-dimensional no-capture property of the center piece
makes their proportion $o(1)$: its initial cutoff is $B^{1+o(1)}$,
well above every bounded degree. Average both conclusions with weights
$n_Y/|P_y|$. They remove only $o(1)$ of the center--component incidence
mass.

At every retained incidence, the point-projection formula gives
\begin{equation}\label{dep:projection-formula}
 t_p(Y)\deg\overline{\pi_pY}
       =\deg Y-\operatorname{mult}_pY
       \ge\tfrac12\deg Y,
\end{equation}
where $t_p(Y)$ is the generic map degree. For a fixed $p$, group
components having the same projected image $Z$ and define their depth
to be the sum of their map degrees. If all retained groups have depth
at most $R\exp(\eta s)$, then any subcollection of degree $a$ has
projected degree at least
\begin{equation}\label{dep:depth-to-degree}
 a\big/\bigl(2R\exp(\eta s)\bigr).
\end{equation}
Thus it suffices to bound depth after negligible weighted incidence
deletions. Fixed constants and the errors in
\eqref{dep:carrier-budget} are absorbed by choosing
$\eta>0$ sufficiently smaller than the desired $\epsilon$.

Whenever the deleted incidence mass is $\alpha=o(1)$, Markov's
inequality leaves all but $\sqrt\alpha$ of the centers with deleted
target mass at most $\sqrt\alpha$. Every later subcollection of
fixed positive mass still has positive mass after this deletion.
This is the uniformity mechanism in the proposition. No union bound
over the possible subcollections will be needed.

\paragraph{The first-cut bound.}
Use the system at the first plateau after a drop. Its equations have
degrees at most $D_1\exp(o(s))$, and every current descendant is
contained in this first system. The regularity conclusion of
Proposition~\ref{prof:regularity} makes its relevant carrier unique
and smooth at all but $o(|P_y|)$ assigned target points. For each such
point $q$, centers $p$ in its affine tangent flat form a
$B^{-c_0+o(1)}$ fraction of the source piece.

If, for a center and a current descendant $Y$, the generic joining ray
were contained in the first system, then the same containment holds
for every point of $Y$ where the ray is defined, by polynomial
identity. At an assigned regular point $q$ the line would lie locally
in its unique first carrier, so its direction would be tangent there.
The preceding averaging therefore discards all such
center--descendant incidences at negligible mass. On every other
generic ray at least one first-system equation restricts nontrivially.
Its degree bounds the number of distinct intersections with all
current target components by $D_1\exp(o(s))$. In characteristic zero,
the generic fibers defining the map degrees consist of distinct
points. This proves the depth bound for \eqref{dep:first-R}.
Only regularity of the carrier is used here; the original path
equations need not themselves have full Jacobian rank.

\paragraph{Centers in the narrower window.}
Suppose $v/s\to0$. The center piece can be simultaneously isolated as
zero-dimensional components of equations of degree
\[
 L=B\exp(o(s)),\qquad |P_x|=N\exp(o(s)),
\]
where the second error is signed. Fix $Y\in\mathcal Y$, and choose
$y$ general on it, also general for the projections from all finitely
many actual centers. For each center, count the additional moving
roots $z\ne y,p$ on target components whose images agree with
$\overline{\pi_pY}$. These points are generic on their components.
Here this includes avoiding their intersections with the other members
of the same finite family $\mathcal Y$. Their number is the group
depth minus the one fixed root $y$.

Parametrize such pairs by
\[
 p=y+t(z-y),\qquad z\in\bigcup_{Z\in\mathcal Y}Z.
\]
Impose the center-isolating equations on this parametrization.
At every counted solution the image point $p$ is isolated in the
center system. At its root $z$, the target union is locally just
the one member of $\mathcal Y$ containing that root, by the preceding
genericity condition. The generic ray has finite intersection with
that component, so the fiber is locally finite. Thus the solution is
isolated on the variety parametrized by $(z,t)$.
This variety has dimension $l+1$ and degree at most a dimensional
constant times $\sum_Z\deg Z$; the pulled-back equations have degree
at most $2L$. The isolated-intersection form of B\'ezout therefore
bounds the total number, over all centers, by
\begin{equation}\label{dep:narrow-root-count}
 C e_j\exp(o(s))L^{l+1}.
\end{equation}
Possible larger components of the center equations away from their
isolated sample points do not affect this count.

Divide \eqref{dep:narrow-root-count} by $|P_x|$. Its value, to the
stated accuracy, is
\[
 \frac{e_j B^{l+1}}{N}
 =B\exp\left(s\sum_{i\le j}b_i\right)
 =B\exp\left(-s\sum_{i>j}b_i\right)=R.
\]
The fixed root adds only one; $R=B^{1+o(1)}\to\infty$.
For every fixed logarithmic slack $\eta>0$, Markov's inequality bounds
the proportion of centers whose depth at this $Y$ exceeds
$R\exp(\eta s)$ by $\exp(-\eta s+o(s))$.
Weighted averaging in $Y$, followed by the row deletion already
described, proves \eqref{dep:narrow-R}.

\paragraph{Centers in the wider window: the generic lists.}
Suppose $v/s\to\infty$, $b_{j+1}>0$, and $l<d-1$.
Choose an interior center plateau just to the right of zero, before
the first strictly positive center drop. Its carrier dimension is
\[
 k=\#\{i:b_i>0\}\ge l,
 \qquad k<d.
\]
All componentwise proper cuts of degree $B\exp(Cs)$, with $C$ fixed,
are negligible for the center piece. To verify the degree comparison,
choose a fixed positive interior testing time $\tau$ on this center
plateau. Since $s=o(v)$,
$B\exp(Cs)=o(B\exp(\tau v))$, and leave a second interior time for
the pooling step of Proposition~\ref{prof:regularity}.

For each $p\in P_x$, fix one carrier $X(p)$ of this center plateau
which contains it. The distinct values form a source family
$\mathcal X$, with disjoint assigned sets
$S_X=\{p\in P_x:X(p)=X\}$. Apply the nonreal-flat deletion in
Proposition~\ref{prof:regularity} to this source assignment and to
the fixed target assignment on $\mathcal Y$. Use it to omit every
nonreal flat member from both input families, and omit any member
of zero assigned mass. The removed assignments have source mass
$r_x=o(|P_x|)$ and target mass $r_y=o(|P_y|)$. Discard every incidence whose center belongs to the
removed source set or whose target component is removed. Their
total normalized weight is at most
\[
 \frac{r_x}{|P_x|}+\frac{r_y}{|P_y|}=o(1).
\]
Call the surviving distinct families $\mathcal X_0,\mathcal Y_0$,
keeping their assigned sets, and put
\[
 n_x=\sum_{X\in\mathcal X_0}|S_X|=(1-o(1))|P_x|,
 \qquad
 n_y=\sum_{Y\in\mathcal Y_0}n_Y=(1-o(1))|P_y|.
\]
Both are at least $N\exp(-o(\log B))$, and every flat member of
these families is real. When $k=l$, the degree and assigned-mass
budgets, using $v,s=o(\log B)$, give at most $B^{d-l+o(1)}$ flat
members in each family and normalized weight at most
$B^{-(d-l)+o(1)}$ for each one. Thus
Lemma~\ref{geo:generic-depth} applies with error scale $\log B$ and,
when $k=l$, with $e=B^{d-l}$. Its removed product weight has
normalized incidence weight
\[
 \frac{n_xn_y}{|P_x||P_y|}\,o(1)=o(1)
\]
in the incidence normalization fixed above.

The lemma supplies a list $\mathcal Y_X\subseteq\mathcal Y_0$ for
each $X\in\mathcal X_0$ such that for every $Y\in\mathcal Y_X$,
a general line through $x\in X,y\in Y$ has
at most
\[
 q=B^{1/2}
\]
moving roots on that list. Using a smaller fixed exponent in the
lemma, if necessary, ensures this displayed bound for all sufficiently
large $B$. This is a bound on all possible generic list meetings,
not just on the roots visible in the finite sample. For a fixed
general $x$, a proper subvariety of an $l$-dimensional listed target
has projected dimension less than $l$, so cannot supply roots over
the generic direction image of $Y$.
For a retained actual center $p$, use only the list indexed by its
fixed carrier $X(p)$, further restricted by the earlier incidence
deletions at $p$.

We prove that excessive depth at an actual center is detected by a
proper cut on $X$ of degree $B\exp(O(s))$. Put
\[
 r=\left\lceil B\exp(-2s)\right\rceil.
\]
Because $s=o(\log B)$, we have $q=o(r)$ and $r=B^{1+o(1)}$.
Fix $X$ and $Y\in\mathcal Y_X$ for the following construction.

The generic list bound does not yet control actual centers.  We
encode intersections with a joining line by two binary forms.
Excessive retained depth will force a large rank decrease in their
Sylvester matrix.  Multiplicity compression will then place these
centers in a proper cut of degree fitting the center plateau.

\paragraph{A fixed generic target point and two binary forms.}
Write the current target path locus as
\[
 Z=V(f_1,\ldots,f_m),\qquad \deg f_i\le L_0=B\exp(O(s)).
\]
This is the fixed full path locus, including components omitted from
the generic-depth input family.
Every retained target is an irreducible component of $Z$, as proved
above. Choose $y\in Y$ general with two sets of conditions.
First, the generic list-root bound and the non-tangency conditions
must hold for generic $x\in X$. A dense open subset of $X\times Y$
has a dense open fiber over general $y$, so this is possible.
Second, for each of the finitely many retained actual centers $p$
assigned to $X$ and each listed target $Y'$ with $l$-dimensional image equal to
$\overline{\pi_pY}$, require all points of the generic fiber over
$\pi_py$ to be affine, distinct, different from $p$, and outside every
other irreducible component of $Z$.
The generically finite map from $Y'$ has a dense open image locus
with these properties: every excluded proper subset of $Y'$ has
image of dimension at most $l-1$. Pulling back these open sets to
$Y$ and intersecting the finite list preserves a nonempty open set.
Also avoid the finitely many actual centers. This fixes one $y$ for
all the specializations below.

Choose an integer $L\ge\max(L_0,B)$ with $L=B\exp(O(s))$.
Homogenize each path equation to its own degree and take all its
homogeneous monomial multiples to degree $L$. Their span has
projective base locus exactly $V(f_1^{\mathrm h},\ldots,f_m^{\mathrm h})$:
at every projective point some monomial of each prescribed degree is
nonzero. This base locus may contain extra components at infinity;
we retain all of them. The closure of each retained affine target
is still a component, since any larger component containing it
would meet the affine chart and contradict the component property.

Parametrize the line through $x$ and $y$ by
\[
 \iota_x([U:V])=[U:Uy+V(x-y)].
\]
For $x\ne y$, this identifies $\mathbb P^1$ with the line.
The parameters of $y,x$, and infinity are respectively
$[1:0],[1:1],[0:1]$, independently of $x$.
Every restricted equation is a binary form of degree $L$, whose
coefficients are polynomials of degree at most $L$ in $x$.
The general line is not contained in the base locus. At the chosen
smooth point $y$, outside its other components, such containment
would put its direction in $T_yY$, contradicting the retained
non-tangency condition.

Over $K=\mathbb C(X)$, let $H$ be the greatest common divisor of the
restricted system and put $g=\deg H\le L$.
After division by $H$ the binary linear system has no base point.
Two generic members of it have no common zero: choose one nonzero
member, then choose the second to avoid its finitely many zeros.
Consequently two generic constant complex linear combinations of
the original degree-$L$ forms restrict to binary forms $F_x,G_x$
whose generic gcd is exactly $H$.
The good choices form a nonempty algebraic open set over $K$ and
admit constant complex choices: expanding the coefficients of a
nonzero polynomial over $K$ in a complex-linearly independent basis
shows that it cannot vanish on every tuple of complex constants.
This also covers $g=L$, when the residual forms are constants.

\paragraph{The Sylvester matrix and specialization of its divisor.}
Let $S(x)$ be the $2L$-square matrix of
\[
 K[U,V]_{L-1}\oplus K[U,V]_{L-1}
 \longrightarrow K[U,V]_{2L-1},\qquad
 (a,b)\longmapsto aF_x+bG_x.
\]
Its entries have ambient degree at most $L$ in $x$, and its generic
rank on $X$ is $\rho=2L-g$.
Indeed, after writing $F_x=Hf$ and $G_x=Hh$ with $f,h$ coprime,
the kernel consists of $(hQ,-fQ)$ with $\deg Q=g-1$, and has dimension
$g$ (zero if $g=0$).

Consider an actual center $p$ assigned to $X$ at which the retained
depth over $\overline{\pi_pY}$ exceeds $r$. By the choice of $y$, at least
$r-O(1)$ distinct common roots of $F_p,G_p$ are affine points of
retained components, different from $p,y$, and each lies on no other
component of the fixed base locus. Call these distinguished roots.

The generic gcd divisor specializes with its full length, even when
$X$ is singular at $p$. To see this explicitly, choose an irreducible
algebraic curve through $p$ whose general point lies in the open set
where the preceding generic properties, including the gcd degree,
hold. To obtain it, take successive general linear sections through
$p$ until the local dimension is one, avoiding every
positive-dimensional component of the excluded closed set at each
step. That closed set then has local dimension at most zero.
Choose an irreducible curve component through $p$; it meets the good
open set. Normalize it and use the discrete
valuation ring at a point above $p$. Over its fraction field choose
the gcd form, and scale it so its coefficients are integral and at
least one is a unit. Its reduction $H_0$ is a nonzero homogeneous
form of degree $g$. It divides both specialized forms. In fact, the
quotients are integral: the content law for polynomials over a
discrete valuation ring says that a primitive factor times a quotient
of negative coefficient valuation cannot have all coefficients
integral. Reducing the integral factorizations proves divisibility.
Thus $H_0$ defines a common limiting divisor $D_0$ of length $g$ on
$\mathbb P^1$.

At most $q$ distinguished roots lie in the support of $D_0$.
For this assertion, pass to a finite extension splitting the generic
gcd and a valuation ring above the chosen one. Each distinct generic
root has a unique projective specialization, obtained by scaling its
two coordinates to be integral with one a unit. Each physical root
belongs to some fixed irreducible component of the projective base
locus. Closedness forces its specialization to remain in that
component. A distinguished point lies only on its own retained
component, so a branch specializing to it must already have come
from that retained component. It cannot be either immobile endpoint,
whose fixed parameters specialize to $y$ and $p$, nor a root fixed
at infinity. The generic list bound permits at most $q$ remaining
distinct roots. Their specializations have at most $q$ support points,
regardless of the multiplicities carried by the gcd divisor.

If $F_p,G_p$ are not both zero, their common divisor therefore has
degree at least
\[
 g+r-O(1)-q.
\]
This includes the case that exactly one form is zero: its common
divisor with the other is the other degree-$L$ form, and the matrix
has rank $L$. In all these cases
\[
 \rho-\operatorname{rank}S(p)\ge r-q-O(1)\ge r/2
\]
for large $B$. If both forms vanish, $S(p)=0$ and the rank decrease
is $\rho\ge L\ge r$, giving the same conclusion.
Working throughout with binary forms of degree $L$ records any loss
of an affine leading coefficient as a root at infinity; hence that
phenomenon introduces no exception to the comparison.

We have shown that depth exceeding $r$ forces a matrix rank decrease
of at least $r/2$.  It remains to turn this into a proper polynomial
cut of degree $B\exp(O(s))$, to which center no-capture applies.

\paragraph{Multiplicity compression and the center no-capture bound.}
Choose a $\rho$-square minor $Q(x)$ nonzero generically on $X$.
It has degree at most $2L^2$ and has ambient multiplicity at least
$r/2$ at every excessive actual center. For clarity, at such a center
the selected submatrix has rank at most $\rho-r/2$. Constant
invertible row and column operations put its value into rank normal
form. At least $\lceil r/2\rceil$ rows then have every entry in the ambient
maximal ideal at $p$, so every determinant term belongs to its
$\lceil r/2\rceil$-th power.
No smoothness of $X$, or agreement between ambient generic rank and
generic rank on $X$, is needed for this conclusion.

Choose a test point of $X$ at which $Q$ is nonzero. The multiplicity
compression lemma, Lemma~\ref{poly:compression}, produces a polynomial
which vanishes at all these excessive centers, is nonzero at the
test point, and has degree
\[
 O(L^2/r)=B\exp(O(s)).
\]
It is consequently a proper cut on $X$.
For each fixed $Y$, these cuts can be chosen separately on the center
components $X$ with $Y\in\mathcal Y_X$. Their degree bound is uniform.
The componentwise pooling and no-capture statement on the chosen
center plateau makes the captured center mass $o(|P_x|)$.
This is uniform in the choice of $Y$: a contrary sequence of targets
would contradict the same no-capture statement, which allows arbitrary
polynomials of the specified degree. Averaging with the normalized
weights $n_Y/|P_y|$ therefore removes only $o(1)$ of all incidences.

On the remaining incidences, every group in the retained part of
$\mathcal Y_{X(p)}$ at $p$ has depth at most $r$.
Deleting further targets cannot increase depth.
Apply the row deletion and \eqref{dep:depth-to-degree}. Since rounding
$r$ costs a factor $1+o(1)$, this proves \eqref{dep:wide-R} and
completes the proof of the proposition.
\end{proof}

\begin{remark}\label{dep:uniformity-remark}
The proof bounds incidence mass before discarding centers. It does
not claim a uniform depth estimate for every individual original
target at every center. The conclusion needed later is that, after
one fixed deletion, every target subcollection
of fixed positive row mass satisfies the projected-degree bound.
This permits adaptive choices such as the family on which later
regularity equations fail. It also explains why no bound on the
number of target components enters a union bound.
\end{remark}
